% !TeX root = ../组合数学笔记.tex
% 由主文档通过 \input 载入；请编译主文档。
\setcounter{chapter}{0} % 保留原书章号。
\chapter{绪论、组合对象与证明方法}
\label{ch:objects}

从四个人中选两名代表有 $6$ 种方案，选一名正代表和一名副代表却有 $12$ 种方案。差别不在计算技巧，而在计数对象：前者是二元素子集，后者是由两个不同元素组成的有序对。组合问题的第一步因此不是选择公式，而是确定什么算作一个对象、哪些对象合法、何时两个对象相同。

\section{计数、存在、构造与优化}

设 $X$ 是满足给定限制的有限对象集合。\emph{计数问题（enumeration problem）}要求 $|X|$，存在性问题要求判断 $X$ 是否非空，构造问题要求实际给出 $X$ 的一个元素；若另给代价函数 $w:X\to\mathbb R$，优化问题要求在 $X$ 中使 $w$ 最小或最大。后者只有在可行集非空时才有可行解；有限性保证非空可行集上的实值函数取得极值，却不保证容易找到极值点。

例如，每个学生有若干可选课题，且每个课题至多分给一人。可以问合法分配的数量、是否存在覆盖所有学生的分配、怎样找到分配，或怎样使总偏好得分最大。同一情境产生四个不同任务。给出一个方案只能证明存在，检查许多失败方案一般不能证明不存在；一个计数公式也不自动附带有效的构造算法。

\section{集合、映射与组合编码}

\begin{definition}
映射 $f:X\to Y$ 若满足 $f(x)=f(x')\Rightarrow x=x'$，称为\emph{单射（injection）}；若每个 $y\in Y$ 都有原像，称为\emph{满射（surjection）}；同时为单射和满射时称为\emph{双射（bijection）}。集合上的关系若具有自反性、对称性、传递性，称为\emph{等价关系（equivalence relation）}。
\end{definition}

对有限集合，单射给出 $|X|\le|Y|$，满射给出 $|X|\ge|Y|$，双射给出 $|X|=|Y|$。等价关系把集合划分成互不相交的等价类；将某些变换后的对象视为同一个方案，就是在计数等价类。构造编码时必须记录解码所需的信息，否则只能得到一个多对一映射，不能直接断言等数。

\begin{example}[子集与零一串]
对子集 $A\subseteq[n]$，令第 $i$ 位为 $1$ 当且仅当 $i\in A$，得到一个长度为 $n$ 的零一串。反向取出所有值为 $1$ 的位置，便恢复 $A$；两步互逆。因此 $[n]$ 的子集与长度为 $n$ 的零一串一一对应。此编码还保持“子集大小”等于“串中 $1$ 的个数”，可以直接处理附加限制。
\end{example}

\section{加法、乘法与等重计数原理}

\begin{theorem}[基本计数原理]
若有限集合 $X$ 是两两不交集合 $X_1,\ldots,X_m$ 的并，则 $|X|=\sum_i|X_i|$。若一个构造分成 $r$ 步，对任意合法的前 $i-1$ 步选择，第 $i$ 步都有恰好 $a_i$ 种延续，且最终选择序列与对象一一对应，则对象数为 $\prod_i a_i$。若 $f:X\to Y$ 满足每个 $y\in Y$ 恰有 $d>0$ 个原像，则 $|X|=d|Y|$。
\end{theorem}
\begin{proof}
第一式中每个元素恰在一项中被计数。第二式对步数归纳：每个长度为 $i-1$ 的合法前缀有 $a_i$ 个不同扩展，所以前缀总数逐步乘以 $a_i$。第三式将 $X$ 分成各纤维 $f^{-1}(\{y\})$，对这些互不相交的 $d$ 元集合应用第一式。
\end{proof}

前两条分别称为\emph{加法原理（sum rule）}和\emph{乘法原理（product rule）}。乘法原理不要求各步选择的具体内容互不影响，但要求相应延续数恒定；若第一步选 $x\in A$ 后有 $b(x)$ 种延续，总数应为 $\sum_{x\in A}b(x)$。例如先选一个数 $i\in[n]$，再选 $j\in[i]$，得到 $n(n+1)/2$ 个有序对，而非 $n^2$ 个。

等重条件是除法的依据。从四人中顺次选两人有 $4\cdot3$ 种，每个二元素子集恰有两种排列，故子集数为 $12/2=6$。若某些对象只有一种编码、另一些有两种编码，则不存在一个可以统一除去的“重复次数”。此外，空串与空排列各是一个对象，所以零步构造的方案数是 $1$；这与空集的元素数为 $0$ 并不矛盾。

\section{归纳、双射与最小反例}

\emph{数学归纳法（mathematical induction）}要求验证起点，并由 $P(n)$ 推出 $P(n+1)$。\emph{强归纳法（strong induction）}允许在证明 $P(n)$ 时使用所有更小且不低于起点的情形。二者可以由非负整数的良序性质解释：若存在反例，最小反例的所有较小情形成立，归纳步骤便排除了它。因此使用递归分解时，必须说明参数确实减小，并覆盖最小对象。

\begin{proposition}
对 $n\ge0$，$|2^{[n]}|=2^n$；对 $n\ge1$，偶数大小子集和奇数大小子集各有 $2^{n-1}$ 个。
\end{proposition}
\begin{proof}
子集的每个元素有选或不选两种状态，故由乘法原理得到第一式；也可按是否含 $n$ 分成两个与 $2^{[n-1]}$ 等大的类，以初值 $|2^{\varnothing}|=1$ 归纳。对第二式，固定元素 $1$，把包含 $1$ 的子集删去 $1$，把不包含 $1$ 的子集加入 $1$。此操作改变大小的奇偶性，连续做两次恢复原集合，因此给出两类之间的双射。二者和为 $2^n$，各占一半。
\end{proof}

该证明展示了两种组织方式：递推按较小对象分解，双射直接匹配两类对象。选择方法的标准是能否清楚解释对象，不是公式看起来是否简短。

\section{习题}
\begin{enumerate}
  \item\label{ex:objects-1} 从五人中选三名代表，再从代表中指定一名召集人，求方案数，并用另一种选择次序核对。
  \item\label{ex:objects-2} 设有限映射 $f:X\to Y$ 的每个纤维至多含 $d$ 个元素，证明 $|X|\le d|Y|$；给出等号成立的条件。
  \item\label{ex:objects-3} 说明从 $[n]$ 到空集的映射数，以及从空集到任意集合的映射数。区分 $n=0$ 与 $n>0$。
  \item\label{ex:objects-4} 用归纳法证明：有限集合 $A$ 的所有子集的大小之和为 $|A|2^{|A|-1}$（$|A|\ge1$），再尝试给出双重计数证明。
\end{enumerate}

